\documentclass[11pt]{article}
\usepackage[a4paper,margin=25mm]{geometry}
\usepackage{amsmath,amssymb,amsthm,booktabs,longtable,hyperref}
\newtheorem{theorem}{Theorem}
\newtheorem{lemma}[theorem]{Lemma}
\newcommand{\Sh}{\operatorname{Sh}}
\title{Shadows and residues in finite coset geometries}
\author{Human Compression project: private partial recovery}
\date{}
\begin{document}\maketitle
\noindent The purpose of this note is to decide geometric shadow intersection
properties using the same finite coset actions in every residue. The orbit
and witness reductions are those of Jacobs and Leemans [JL, Section 3],
restated in the independently reviewed working note. The extension below
makes the residue quantifiers and the input geometry tests explicit. It
reuses the previously accepted flag-constrained covering proof and the
panel-recognition arguments; it is not a new claim to the source algorithm.
\section{Orbit and witness recognition}
An incidence geometry has a finite set of types $I$, reflexive symmetric
incidence, and no incident distinct objects of the same type. Every flag
extends to a chamber. Suppose a finite group $G$ acts transitively on flags
of each type. Fix a chamber $C=\{p_i:i\in I\}$, use right actions, and write
\[
 H_i=G_{p_i},\quad H_T=\bigcap_{t\in T}H_t,\quad H_\varnothing=G,
 \qquad C_T=\{p_t:t\in T\}.
\]
The $i$-shadow $\Sh_i(F)$ is the set of type-$i$ objects incident with every
member of $F$. This geometric shadow property is distinct from the
subgroup intersection property for distinguished generators of a C-group.

\begin{lemma}\label{orbit}
For every $i,T,h$ we have
\[
             \Sh_i(C_T^h)=(p_i^{H_T})^h.
\]
\end{lemma}
\begin{proof}
If $i\in T$, both sides are the singleton $\{p_i^h\}$.
Otherwise a type-$i$ object $y$ incident with $C_T$ completes it to a flag
of type $T\cup\{i\}$. Flag transitivity gives an element carrying
$C_{T\cup\{i\}}$ to that flag while fixing each member of $C_T$.
It belongs to $H_T$ and carries $p_i$ to $y$. Conversely every element of
$H_T$ fixes $C_T$, so its image of $p_i$ is incident with $C_T$.
Apply $h$ to obtain the stated equality.
\end{proof}

The objects of type $i$ are identified with right cosets $H_i\backslash G$.
Thus the lemma uses one permutation action for each type. All shadows in
that action are orbits of intersections of the chamber stabilizers.
The recognition theorem in the current common theory for finite geometries
can establish the required flag transitivity: for every $i\notin T$, the
intersection $\bigcap_{t\in T}p_i^{H_t}$ must equal $p_i^{H_T}$.
Induction on $|T|$ proves that this condition is sufficient, starting with
the empty flag. It checks the entire set of eligible extensions, not a
sample of chambers.

\begin{theorem}\label{shadowtest}
Fix a type $i$. The following condition is equivalent to the geometric
shadow intersection condition at type $i$: for every element $x$ and flag
$F$, the intersection $\Sh_i(x)\cap\Sh_i(F)$ is empty or equals
$\Sh_i(F')$ for a flag $F'$ incident with both $x$ and $F$.

For each $j\ne i$, put $K=H_i\cap H_j$. For every nonempty
$T\subseteq I\setminus\{i\}$, choose representatives $h$ of
\[
                  (H_i\cap H_T)\backslash H_i/K.
\]
For each such representative form
\[
                  X=p_i^{H_j}\cap(p_i^{H_T})^h.
\]
If $|X|>1$, there must be a subset $U\subseteq I\setminus\{i\}$
and an element $q\in K$ such that
\[
 (p_i^{H_U})^q=X,
 \qquad p_u^q\in(p_u^{H_T})^h\quad\hbox{for every }u\in U.
\]
For a fixed $U$ it suffices to take one representative of each right
coset of $K\cap H_U$ in $K$.
\end{theorem}
\begin{proof}
Cases where $x$ has type $i$, $F$ contains type $i$, or $F$ is empty have
an immediate witness whenever the intersection is nonempty. Otherwise
choose $p$ in the nonempty intersection. The pair $\{p,x\}$ is a flag of
types $\{i,j\}$; normalize it to $\{p_i,p_j\}$. The normalized $F$ lies
in the residue of $p_i$. Flag transitivity on $T\cup\{i\}$ expresses it
as $C_T^h$ with $h\in H_i$. Its stabilizer inside $H_i$ is $H_i\cap H_T$.
The remaining choices of normalization act by $K$, giving precisely the
displayed double cosets. Lemma~\ref{orbit} identifies its intersection
of shadows with $X$.

The singleton case has the witness consisting of its unique type-$i$
object. If $|X|>1$, a witness cannot contain type $i$. Let its type be
$U$. Because $p_i\in X$, it is incident with $p_i$, and by assumption it
is incident with $p_j$. Hence its union with $\{p_i,p_j\}$ is a flag.
Flag transitivity expresses the witness as $C_U^q$ for $q\in K$.
Two elements of $K$ give the same such flag exactly when they are in
the same right coset of $K\cap H_U$. Lemma~\ref{orbit} gives the first
displayed condition. Incidence with $C_T^h$ says that, for every $u\in U$,
its type-$u$ object belongs to $\Sh_u(C_T^h)$, giving the second.

Conversely $C_U^q$ is a flag. Since $q$ fixes $p_i$ and $p_j$, it is
incident with those two objects. The second displayed condition makes
it incident with $C_T^h$, and the first identifies its $i$-shadow with
$X$. Transporting it back gives the required original witness. Every
normalization and every orbit representative preserves these conditions,
which proves sufficiency as well as necessity.
\end{proof}

There are two useful exact shortcuts. The candidate type $U$ can be
discarded unless $|p_i^{H_U}|=|X|$. Also, if $p_j$ is incident with
$C_T^h$ and $|X|$ is the smaller of the two original shadow sizes, one
of the original flags is already a witness. Both are consequences of
the theorem and preserve its complete search space.


\section{Every residue and every rank}
For $J\subseteq I$, let $R_J$ be the residue of the base flag $C_J$.
Its types are $I\setminus J$. Empty residues are not tested for an
existential type condition.
\begin{theorem}\label{residue}
Suppose the coset system is flag transitive. The map
\[
 (H_J\cap H_i)h\longmapsto H_i h,\qquad h\in H_J,
 \quad i\in I\setminus J,
\]
is an isomorphism of the coset geometry inside $H_J$ with $R_J$.
In particular $H_J$ is flag transitive on $R_J$.
For each $k\in\{2,\ldots,|I|\}$ it suffices to test precisely the base
residues $R_J$ with $|I\setminus J|=k$. On each of them apply
Theorem~\ref{shadowtest} with group $H_J$ and parabolics $H_J\cap H_i$.
All type tests in all these residues give $(\mathrm{IP})_k$;
at least one type test in each residue gives $(\mathrm{WIP})_k$.
The conjunction over $k=2,\ldots,|I|$ gives respectively IP and WIP.
\end{theorem}
\begin{proof}
An object of type $i$ incident with $C_J$ extends that flag. Flag
transitivity therefore writes it as $H_i h$ for $h\in H_J$. Two such
representatives give the same object exactly when they give the same
right coset of $H_J\cap H_i$ inside $H_J$. This proves object bijectivity.
If the small cosets intersect, their common representative gives incidence
of their images. Conversely, two incident images together with $C_J$
form a flag. Flag transitivity gives a representative in $H_J$ for both
images, so their small cosets intersect. The same argument, applied to
an arbitrary flag of the residue together with $C_J$, proves flag
transitivity under $H_J$.

Every flag of type $J$ is a translate of $C_J$, so all its residues are
isomorphic with their type labels retained. This proves that one base
residue per subset is sufficient; it does not identify different subsets
or assume their residues are isomorphic. Inside a fixed residue the
preceding theorem is an equivalence at each remaining type. Universal
quantification of those types is IP at that rank; the existential type
for WIP is chosen separately in each residue. These are precisely
[JL, Lemmas 2.1--2.2]. Rank-one shadow intersections are empty, singletons,
or the whole sole type, with the corresponding immediate witnesses.
Rank zero has no available type and is excluded, as in those lemmas.
\end{proof}
There are $[G:H_J]$ flags of type $J$, so the report records that many
actual residues represented by each base residue. It computes every
remaining type, even when one failed type would already refute IP.
Thus its WIP output does not silently use a fixed global type.

\section{Geometry hypotheses are separate checks}
The supplementary rows are not accepted because their total matches a
published table. They define coset incidence structures, which need not
be flag transitive. We first check the orbit equality preceding
Theorem~\ref{shadowtest} for all $i,T$. This condition is necessary by
Lemma~\ref{orbit}. For sufficiency, normalize an existing flag of type
$T$ inductively. Every object eligible to extend it belongs to the
intersection of the singleton shadows, hence to the $H_T$ orbit by the
checked equality. It follows that each extended flag is a translate of
the corresponding base flag. Induction covers every flag type, and
in particular every flag extends to a chamber.

Write $K=H_J$ and $P_i=K\cap H_i$ for a residue. The following tests
are the actual source hypotheses, with kernels of actions allowed.
\begin{lemma}\label{filters}
For a flag-transitive coset geometry the following conditions are exact.
\begin{enumerate}
\item Firmness holds if and only if $[H_{I\setminus\{i\}}:H_I]\ge2$
for every $i$.
\item Residual connectedness holds if and only if
$\langle P_i:i\in I\setminus J\rangle=K$ for every residue of rank
at least two.
\item Residual weak primitivity holds if and only if in each residue of
positive rank at least one transitive action on $P_i\backslash K$ is
primitive of degree at least two.
\item The local condition $(2T)_1$ holds if and only if, in every
rank-one residue, its point stabilizer has exactly two orbits on its
objects and there are at least two objects.
\end{enumerate}
\end{lemma}
\begin{proof}
The objects completing the base panel missing type $i$ are the cosets
of $H_I$ in $H_{I\setminus\{i\}}$. Flag transitivity covers every panel,
proving the first claim. For the second, a coset path starting at the base
chamber can multiply its representative successively by elements of the
$P_i$. Conversely each such multiplication is realized by a path through
the object fixed by that subgroup. Hence the connected component of the
base chamber is exactly the orbit of their join, and all components are
its cosets. The argument is applied inside each $K$; it is valid for
rank two as well as higher ranks. The third statement is the definition
of weak primitivity applied to each residue. Equivalently $P_i$ must be
maximal proper in $K$: blocks containing the base coset correspond to
intermediate subgroups. This equivalence does not require a faithful
action. Finally a transitive action is two-transitive exactly when the
stabilizer of a point is transitive on the other points. Its two orbits
are then that point and its complement. Again this depends on the
permutation image, not on a faithful ambient action.
\end{proof}
These are the panel and coset-path arguments already used in the accepted
Hall--Janko census [D]. The implementation tests the entire action, not
only subgroup orders or a presumed diagram.

\section{Rank two without a diagram assumption}
In a rank-two incidence structure the only nonautomatic shadow test
compares the neighborhoods of two distinct objects on the same side.
At a fixed shadow type $i$, it passes exactly when every such intersection
has size zero, one, or the whole type-$i$ set. Indeed a witness of type
$i$ has singleton shadow; a witness of the other type cannot be incident
with both distinct objects of that type. The empty flag has the whole
shadow set. These are all possible witness flags. In particular comparing
an object with itself is automatic and is excluded from this pair test.

Both type tests hold if and only if the graph has no four-cycle or is
complete bipartite. If a four-cycle is present, its two vertices on one
side have at least two common neighbors, hence all vertices of the other
side. Every pair on that other side consequently has at least two common
neighbors and must meet the entire first side. This proves completeness.
The converse follows immediately. This is the accepted direct proof
[Cover], recovering the partial-linear-space/generalized-digon criterion
used in [JL, Section 3]. We still record the two directional tests
separately, and do not replace the weak condition by an assumed diagram.

For a coset action, fix the first object of the opposite type $j$.
Representatives of $P_j\backslash K/P_j$ cover its possible second
objects. Discard the identity double coset, which represents the same
object. Compute the two neighborhoods by their actual coset-action
orbits. Their intersections then perform the criterion just proved.

\section{Bounded archive application}
The parser reads only literal permutations (and the literal identity)
from the supplementary archive. It removes comments before applying the
last active reassignment of \texttt{geo}. It records both the construction
index and the resulting active index. Neither is asserted to be a
published geometry label. A commented selection with an empty placeholder
in the rank-four Mathieu file is not executed.

For the Mathieu group $M_{11}$, the current complete checks cover thirteen
active rank-three rows and eleven active rank-four rows. Every row passes
all six source filters: flag transitivity, firmness, residual connectedness,
RWPRI, $(2T)_1$, and $(\mathrm{IP})_2$. All residues of rank at least two
are checked, not only the original geometries.
\begin{center}
\begin{tabular}{ccrrr}\toprule
Input rank & tested residue rank & IP & weak only & not weak\\\midrule
3&2&13&0&0\\3&3&6&1&6\\
4&2&11&0&0\\4&3&2&0&9\\4&4&2&0&9\\\bottomrule
\end{tabular}
\end{center}
The full rank-three IP rows in active archive order are $1,3,9,10,11,13$;
row $2$ has WIP but not IP. In rank four, rows $1,3$ have full IP and
all other rows fail full WIP. The latter are construction rows $10,2$.
These aggregate values agree with the corresponding entries of Tables 1
and 3, but that agreement is a comparison after the computation, never
an input filter or a proof of a published label. The source names the
positive rank-four geometries $4.1,4.4$ in its earlier classification;
that exact identification remains pending.

A failed type test records the actual normalizing permutation, flag types
and intersection. Every possible witness type and its stabilizer cosets
are then exhausted by the checker. Positive decisions exhaust every
required double-coset orbit. The GAP finite-group routines are an explicit
trusted computational input; no Lean formalization of these catalogue
checks is claimed.

\section{The active M12 files and their filters}
The other bounded batch covers all twenty active rank-three rows, all
twenty active rank-four rows, and the single rank-five row of the
$M_{12}$ archive. These are the literal final active lists, not the
numbers assumed from the printed tables. Their groups are explicitly
conjugate to the named natural Mathieu group; the computed conjugating
permutations are retained with the input binding.

Every row passes flag transitivity, firmness, residual connectedness,
RWPRI and $(\mathrm{IP})_2$. Local two-transitivity rejects the
rank-three active rows $15,17,20$ (construction rows $27,29,33$), and
rank-four active rows $4,7$ (construction rows $25,23$).
For rank-three row $17$, the relevant panel action has degree $11$
and order $55$: its point stabilizer has orbit lengths $1,5,5$.
The other failures have degree-three cyclic panel images, whose point
stabilizers have three singleton orbits. These explicit failures do not
depend on interpreting an archived row as a named published geometry.
The two-transitivity test is for the specified $G$ action; no claim
about a larger full automorphism group repairs it.

After these actual filters, the computed counts are as follows.
\begin{center}
\begin{tabular}{ccrrrr}\toprule
Input rank & residue rank & eligible rows & IP & weak only & not weak\\\midrule
3&2&17&17&0&0\\3&3&17&7&1&9\\
4&2&18&18&0&0\\4&3&18&4&0&14\\4&4&18&3&0&15\\
5&2&1&1&0&0\\5&3&1&1&0&0\\5&4&1&1&0&0\\5&5&1&0&0&1\\\bottomrule
\end{tabular}
\end{center}
In particular the rank-five example satisfies every proper-residue
condition, but fails WIP in its full rank. This is why a computation of
proper residues alone cannot certify the full property.

The source's rank-three summary has eighteen rows, with counts
$7,1,10$; the specified active archive under its stated action filters
has seventeen, with counts $7,1,9$. This discrepancy is retained.
The earlier phase-six archive comparison independently flags exactly
these same three rank-three and two rank-four inputs as non-$2T$.
Its independent broader census also records an omitted non-$2T$
rank-three class; adding that class would not repair this filtered
seventeen-versus-eighteen difference. We reuse that prior evidence
without repeating its subgroup census or claiming a newly discovered
literature correction. The rank-four and rank-five aggregate shadow
columns agree with the printed source, but exact published row labels
still require separate identification.

\section{Attribution and remaining boundary}
The supplementary lists and orbit algorithm originate in [JL]. The
present application does not reprove completeness of the source's
sporadic geometry classifications. The entire source contract, the remaining entries of Tables
1--8, their exact published labels, Section 5 on Suzuki and inductively
minimal geometries, and its complete software correspondence remain
pending. In particular active archive sizes need not equal the published
filtered numbers: the two Mathieu $M_{12}$ files each select twenty rows (with the exclusions just proved),
and the supplied $J_2$ rank-four file has one row although the table has
three. Our earlier accepted full Hall--Janko census provides the three
rank-four witnesses independently; its large census need not be rerun.
The executable Magma listing in source Section 6 also requires a matched
software boundary before any numerical comparison; it is not silently
counted only on the source side. No whole-original, question or
numerical-compression credit is claimed.

\begin{thebibliography}{9}
\bibitem[JL]{JL} P. Jacobs and D. Leemans, \emph{An algorithmic analysis
of the intersection property}, LMS J. Comput. Math. 7 (2004), 284--299,
fixed primary PDF and its supplied supplementary archive.
\bibitem[Cover]{Cover} Human Compression project, \emph{Geometric
intersection as a flag-constrained covering problem}, accepted
\texttt{phase4\_shadow\_ip.md} and distinct audit. The finite covering
proof is general; its original audit does not certify this new census.
\bibitem[D]{D} Human Compression project, \emph{Stabilizers, extensions and the enumeration of finite geometries}, accepted Topic D common theory, panel
recognition and Hall--Janko census. Source group catalogues and the
accepted finite census remain explicit retained inputs.
\end{thebibliography}
\end{document}
